\section{Evaluation}
\label{sec:eval}

The recorded evaluation addresses three questions.
\textbf{RQ1}: Does \sys fuse dependent reductions that are used in practice?
\textbf{RQ2}: Which states does it introduce?
\textbf{RQ3}: What does reducer derivation cost?

The coverage, selected states and derivation timings below are recorded in synthesis snapshot \code{6c46fb1b02}.
Controlled comparisons of search guidance and numerical accuracy are described separately below; their measurements under the updated protocol remain pending.
\placeholder{Rerun coverage, state selection, derivation timings and numerical acceptance on the current implementation using the updated MPFR protocol; collect the controlled synthesis comparisons described below.}


We selected \nShowcase operators from applications that differ in how their reductions depend on each other (\cref{tab:showcase}).
Their applications include language models (attention, cross-entropy, LayerNorm and its backward pass), reinforcement learning and multimodal pretraining, optimal transport, statistics, wireless communication, numerical linear algebra and statistical physics.
For most of them, libraries already implement the dependent reductions in one pass by hand; the table identifies recorded implementations.
Each operator is written as its definition in TE, together with a hand-written reducer for comparison.
They come from a registry of \nRegistry operators, many of which are variants of these; Appendix~B of the supplementary material reports the fused registry operators, and Appendix~C additional fusion outcomes and reasons for failures.


All programs are float32 with symbolic extents.
The recorded experiments used Python \pythonVersion, SymPy \sympyVersion, Z3 \zThreeVersion and cvc5 \cvcFiveVersion.
The reported derivation times are medians of five runs after a warm-up, on an Intel Xeon Platinum 8468H.
% src: data/fused_ops.json meta (cpu, extent, seed, offset); data/proofs_timing.json meta.repeats

\begin{table}[t]
\caption{Recorded coverage and state discovery. ``Hand-fused in'' lists identified library implementations; a dash means none is recorded. Passes count reduction traversals, excluding elementwise output reconstruction and producers outside the chain. Derived components and selection routes are explained in the text.}
\label{tab:showcase}
\footnotesize
\setlength{\tabcolsep}{3pt}
\begin{tabular}{@{}>{\raggedright\arraybackslash}p{0.24\linewidth}>{\raggedright\arraybackslash}p{0.28\linewidth}ccc>{\raggedright\arraybackslash}p{0.12\linewidth}@{}}
\toprule
\textbf{Operator} & \textbf{Hand-fused in} & \textbf{Passes} & \textbf{States} & \shortstack{\textbf{Derived}\\\textbf{components}} & \shortstack[l]{\textbf{Selection}\\\textbf{route}}\\
\midrule
\tabShowcaseRows
\bottomrule
\end{tabular}
\end{table}

\subsection{Fusion Coverage and State Selection}
\label{sec:eval-coverage}

In the recorded snapshot, \sys fused \nShowFused of the \nShowcase operators.
Their \showChains chains require \showPassesBefore reduction traversals before fusion and \showPassesAfter afterwards, once per chain; the contrastive loss and the Sinkhorn iteration each have two chains, along the rows and along the columns.
The counts concern the reductions being fused; output elementwise computations and producers outside those chains remain separate.
The hand-written reducers use the same number of reduction traversals.
The derived reducers include the online softmax, FlashAttention's normalised reducer, and pairwise centred-moment joins inside LayerNorm and the Pearson correlation~\citep{chan}.
The heat-capacity reduction uses a weighted squared deviation from the mean, whose body falls outside the stated conditions compared in \cref{thm:containment}.


\paragraph{Derived representation.}
\label{sec:eval-states}

The \showChains reducers have \showStates states, including \showDerived derived components in \showDerivedOps operators (\cref{tab:showcase}).
These components are absent from the source reduction results and include replacements introduced by domain repair.
The initial tuple alone suffices for \showOriginalOps operators, augmentation with all correction coefficients for \showRebaseOps, and native synthesis for \showSygusOps.
These are the selection routes reported in the table; the synthesis route includes independently verified enumeration.
For PAPR, the source retains a total, a centred peak and a squared-deviation sum.
Synthesis adds a count and the sum of the input maximum and minimum, instead of storing the two extrema separately.
This proposal comes from native enumeration after an unknown quantified-synthesis result and passes the independent recovery and join checks.
The Householder-type norm selects a candidate expression from context-correction analysis through SyGuS with the injectivity specification, and uses guarded recovery.
Kurtosis and normalised policy loss first repair normalised source states that are undefined on singleton parts.
Other derived states include coefficients of context corrections; Ledoit--Wolf covariance and scalar OLS HC0 illustrate corrections with several such coefficients.
Table~\ref{tab:showcase} reports the showcase representations; Appendix~B includes label smoothing and the other registry operators.
% src: data/fused_ops.json showcase chains[].states, auxiliaries, actions
The accepted tuples fit the default state-count limit supplied by the existence construction.

\subsection{Derivation Cost and Planned Comparisons}
\label{sec:eval-cost}

In the recorded snapshot, reference-state construction took at most \showJudgeMaxMs{}\,ms per operator.
Selection, which derives the join and discharges every proof obligation, took a median of \showSelectMedianMs{}\,ms and at most \showSelectMaxMs{}\,ms over the accepted showcase operators.
\showInvCanonical join invariants are proved by equality after symbolic canonicalisation.
This count describes the final invariants; selection additionally invokes recovery and domain proofs, including SMT, so it does not measure the entire verification workload.
Complete fusion time includes reconstruction, LLVM compilation and numerical acceptance; it must therefore be measured with the MPFR checks enabled.

\paragraph{Planned search comparisons.}
The recorded selection routes show which mechanisms supplied the final states; controlled comparisons are needed to measure their contribution to search.
We will fix the repaired source states, reference tuple, state-count limit and independent checker, and compare a common mixed grammar against reference-only terminals, native enumeration without the synthesis constraint, and single-terminal selection.
The comparisons will report coverage, proposals checked and derivation time under the same resource limits.
They retain the source states and the reference-recovery requirement, isolating candidate discovery, semantic guidance and candidate combination within this formulation.
Appendix~E specifies the common search schedule and measurements.

\paragraph{Planned numerical comparisons.}
\label{sec:eval-accuracy}

The updated numerical evaluation will compare source, hand-written and synthesised programs against the same MPFR result on identical quantised inputs.
Per-output errors and non-finite results will be reported separately from the acceptance decision in \cref{def:gate}.
This comparison measures how closely the emitted programs match the source's accuracy, including cases where the source is more accurate.

To examine the role of state representation, we will additionally compare the source program, direct execution of the actual reference state tuple $\stref$ constructed from the initial source states, and execution of the maintained tuple for variance and softmax.
The two generated versions will retain their emitted reconstruction expressions and use the same compilation settings and reduction order.
Offset and scale sweeps will test the cancellation and range effects that motivate retaining source centring and scaling.
The existing worked examples illustrate these effects; they are not measurements of this compiled three-way comparison.
Separately, we will hold the maintained states fixed and compare the printed context-correction join with the algebraic join, recording their actual emitted expressions.
This isolates the execution formula from state selection; the algebraic join may itself retain local-context expressions.
Appendix~E of the supplementary material specifies the numerical protocol.

\textbf{Performance scope.}\par
\label{sec:eval-perf}

\ifperf
\placeholder{GPU performance of the fused operators against the unfused TE program, PyTorch and the hand-fused libraries of \cref{tab:showcase}; memory traffic and time.}
\else
\ifsubmit\else
\placeholder{GPU performance. If measurements exist before the deadline: time and memory traffic of the \nShowcase operators against the unfused TE program, PyTorch and the libraries of \cref{tab:showcase}. Otherwise this subsection is removed and the claims stay as they are now.}
\fi
\sys does not schedule or generate GPU code yet (\cref{sec:gpu}), so we report no performance numbers.
Each accepted reducer traverses its shared axis once; realised speed also depends on boundary producers, other chains and scheduling.
\fi
